\scanpage{138}
\mspage{-III- 1 -}
\begin{flushright}n° \handwritten{383}\end{flushright}
\begin{center}\S\ 3. \uuline{Groupes engendrés par des réflexions}.\end{center}

\uline{Dans ce} \S, \uline{on désigne par} $E$ \uline{un espace hilbertien réel de dimension finie} $n$, \uline{par} $(x\handwritten{|}y)$ \uline{sa forme bilinéaire}, \uline{et par} $\handwritten{\lVert}x\handwritten{\rVert}$ \uline{sa norme}. \uline{Pour tout hyperplan} (vectoriel ou affine) $H$ \uline{de} $E$, \uline{on note} $s_H$ \uline{la réflexion} (\uline{ou symétrie orthogonale}) \uline{par rapport à} $H$ (cf. \S\ 2, n°2).

1. \uline{Groupes finis engendrés par des réflexions vectorielles}.

\uline{Dans ce} n°, $G$ \uline{désigne un sous-groupe fini du groupe orthogonal de} $E$, \uline{engendré par des réflexions} (par rapport à des hyperplans vectoriels).

Soit $\underline{H}$ un ensemble fini d'hyperplans de $E$ tel que les $s_H$ $(H\in\underline{H})$ engendrent $G$. L'ensemble $W'$ des $x\in E$ invariants pour $G$ est aussi l'ensemble des $x\in E$ invariants par les $s_H$ $(H\in\underline{H})$ ; donc $\displaystyle W'=\handwritten{\bigcap}_{H\in\underline{H}}H$. Le sous-espace $W$ de $E$ orthogonal à $W'$ est stable pour $G$ ; on l'appelle le sous-espace \uline{essentiel} de $E$ (pour $G$). L'application qui à tout $g\in G$ fait correspondre sa restriction à $W$ est un isomorphisme de $G$ sur un sous-groupe fini $G'$ du groupe orthogonal de $W$, engendré par les réflexions $s_{H\handwritten{\cap}W}$ $(H\in\underline{H})$. On dit que $G$ est \uline{essentiel} si $W=E$ ; en tous cas $G'$ est essentiel, ce qui permet de se ramener le plus souvent au cas où $G$ est essentiel.

Nous nous fixerons désormais un ensemble fini $\underline{H}$ d'hyperplans de $E$ tel que

1) les réflexions $s_H$ $(H\in\underline{H})$ engendrent $G$,

2) les hyperplans de $\underline{H}$ sont permutés par $G$.

Notons, à ce propos, que, pour $g\in G$, on a
\[
\begin{array}{l@{\qquad}l}(1)&s_{g(H)}=gs_Hg^{-1}.\end{array}
\]
En principe, toutes les constructions qui suivent dépendent du choix de $\underline{H}$ ; en fait, nous nous abstiendrons de le rappeler dans nos notations et nos énoncés, car le cor.5 du th.1 révélera que l'ensemble des $s_H$ est nécessairement l'ensemble de \uline{toutes} les réflexions de $G$.
% archive sheet number: 137
